% =============================================================================
% Executive Summary
% =============================================================================
\section*{Executive Summary}\label{sec:executive-summary}
\addcontentsline{toc}{section}{Executive Summary}

This paper presents the complete empirical results of the International
Sovereignty Index (ISI) applied to the 27~member states of the European
Union, using vintage-2024 data processed under methodology v1.0.  The
index measures the concentration of bilateral external dependencies
across six structural axes.  Higher scores indicate higher concentration
of counterparty exposure.  The principal findings are:

\begin{itemize}[leftmargin=2em]

\item \textbf{Compressed composite distribution.}
  The ISI composite ranges from \num{0.23567126} (France) to
  \num{0.51748148} (Malta), a span of \num{0.28181022} on the unit
  interval.  The population standard deviation is \num{0.06992188},
  yielding a coefficient of variation of 0.2030.  The interquartile
  range is \num{0.09881618} (P25~= \num{0.29179051}, P75~=
  \num{0.39060669}).  See \cref{tab:summary-stats,fig:histogram}.

\item \textbf{Tier compression.}
  Under the four-tier classification, 23 of 27~countries (85.2\%)
  fall in the \emph{moderately concentrated} tier (ISI $\in [0.25,
  0.50)$).  Only Malta exceeds 0.50 (\emph{highly concentrated});
  three countries (Latvia, Slovakia, France) fall below~0.25
  (\emph{mildly concentrated}).  No country is classified as
  \emph{unconcentrated}.  See \cref{tab:ranking,fig:rank-scatter}.

\item \textbf{Defence and logistics dominate cross-country variance.}
  A covariance-based variance decomposition attributes 35.1\% of
  composite variance to the defence axis and 34.3\% to the logistics
  axis.  Together, these two axes account for 69.4\% of total
  composite variance.  Financial concentration contributes only
  1.0\%.  See \cref{tab:variance-decomp,fig:axis-dominance-bar}.

\item \textbf{Axis dominance is concentrated.}
  In 16 of 27~countries, the defence axis records the highest
  axis-level score; in the remaining 11, logistics is highest.
  No country has its maximum score on the financial, energy,
  technology, or critical-inputs axis.  See \cref{tab:axis-dominance}.

\item \textbf{Three moderate cross-axis correlations.}
  Of 15~pairwise Pearson correlations, three exceed $|r| \geq 0.30$:
  critical inputs--logistics ($r = 0.549$), energy--technology
  ($r = 0.530$), and energy--critical inputs ($r = 0.372$).  The
  remaining 12~pairs are weak ($|r| < 0.30$), indicating that the
  six axes capture largely distinct dimensions.  See
  \cref{tab:corr-pearson,fig:corr-heatmap,fig:corr-network}.

\item \textbf{Robustness to weight perturbation.}
  Under Dirichlet weight perturbation (10\,000 Monte Carlo draws,
  $\alpha_j = 10$), the mean Spearman rank correlation with the
  baseline is $\rho = 0.979$.  The Netherlands exhibits the highest
  rank volatility (standard deviation 3.07 ranks; 90\% band: 10--20).
  See \cref{tab:dirichlet-volatility}.

\item \textbf{Sensitivity to defence axis removal.}
  Leave-one-axis-out analysis yields Spearman correlations ranging
  from $\rho = 0.833$ (defence excluded) to $\rho = 0.993$
  (financial excluded).  Removing the defence axis produces the
  largest rank disruptions; Slovakia moves from rank~26 to rank~27,
  and Malta from rank~1 to rank~2, when defence is excluded.
  See \cref{tab:loo-summary,tab:loo-detail}.

\item \textbf{Z-score standardisation alters rankings materially.}
  When each axis is z-scored before averaging, the rank-order
  correlation with the baseline drops to $\rho = 0.896$, indicating
  moderate sensitivity to the implicit weighting induced by axis
  dispersion.  See \cref{tab:zscore-ranks}.

\item \textbf{Winsorisation preserves full rank order.}
  Capping the defence and logistics axes at their respective 95th
  percentiles and recomputing produces identical rank ordering
  ($\rho = 1.000$), indicating that extreme axis-level values do not
  drive rank inversion.  See \cref{tab:winsorization}.

\item \textbf{Slovakia's zero defence score.}
  Slovakia records a defence axis score of \num{0.00000000}, the
  only zero score in the dataset.  This arises mechanically from
  the absence of bilateral supplier entries for Slovakia in the
  SIPRI arms-transfer data under methodology v1.0.  It does not
  imply the absence of defence procurement.  See
  \cref{sec:country-profiles}.

\end{itemize}
